Cet exercice est un questionnaire à choix multiples. Pour chaque question, une seule des quatre propositions de réponse est exacte. Indiquer sur la copie le numéro de la question et la lettre de la réponse choisie. Aucune justification n’est demandée. Pour chaque question, une réponse fausse, une réponse multiple ou l’absence de réponse ne rapporte ni n’enlève de point.

{\renewcommand{\arraystretch}{1.5}
\begin{tabularx}{\linewidth}{|X|*{4}{>{\centering\arraybackslash}X|}}
	\hline
	\multicolumn{1}{|c|}{\bfseries Question}
	&
	\bfseries Réponse A
	&
	\bfseries Réponse B
	&
	\bfseries Réponse C
	&
	\bfseries Réponse D
	\\
	\hline
	\bfseries Question n°1\newline
	~
	\newline
	\begin{scratch}[scale=0.6]
 	   \blockvariable{mettre \selectmenu{a} à \ovalnum{4}}
 	   \blockvariable{mettre \selectmenu{b} à \ovaloperator{\ovalvariable{a} + \ovalnum{2}}}
 	\end{scratch}
 	~\newline
	&
	~\newline
	À l'issue de ces affectations :\newline
	4 est affecté à a\newline
	2 est affecté à b
	&
	~\newline
	À l'issue de ces affectations :\newline
	4 est affecté à a\newline
	6 est affecté à b
	&
	~\newline
	À l'issue de ces affectations :\newline
	2 est affecté à a\newline
	6 est affecté à b
	&
	~\newline
	À l'issue de ces affectations :\newline
	2 est affecté à a\newline
	2 est affecté à b
	\\
	\hline
	\bfseries Question n°2\newline
	~
	\newline
	\begin{scratch}[scale=0.6]
 	   \blockvariable{mettre \selectmenu{a} à \ovalnum{2}}
 	   \blockrepeat{répéter \ovalnum{4} fois}
     {
     \blockvariable{mettre \selectmenu{a} à \ovaloperator{\ovalvariable{a} + \ovalnum{3}}}
     } 
 	\end{scratch}
 	~\newline
	&
	~\newline
	À l'issue de ces affectations :\newline
	20 est affecté à a
	&
	~\newline
	À l'issue de ces affectations :\newline
	12 est affecté à a
	&
	~\newline
	À l'issue de ces affectations :\newline
	14 est affecté à a
	&
	~\newline
	À l'issue de ces affectations :\newline
	2 est affecté à a
	\\
	\hline
	\textbf{Question n°3}\newline
	Un pavé droit subit une réduction de rapport 0,4 donc :
	&
	~\newline
	son volume est multiplié par $0,4^3$
	&
	~\newline
	son volume est multiplié par $0,4^2$
	&
	~\newline
	son volume est divisé par $0,4^3$
	&
	~\newline
	son volume est divisé par $0,4^2$
	\\
	\hline
	\textbf{Question n°4}\newline
	Un réservoir d'eau a une forme de pavé droit de dimensions exprimées en mètre : \newline
	$1,5\times1,5\times2$\newline
	Son volume en litre est
	&
	~\newline
	450~L
	&
	~\newline
	4~500~L
	&
	~\newline
	4~500~000~L
	&
	~\newline
	4,5~L
	\\
	\hline
\end{tabularx}
}
